% GENERATED FILE. Edit canonical lessons, then run npm run generate:books.
% canonical-source-hash: fnv1a64:799ca515f6070d34
% canonical-lessons: ES-C618-imaginar, ES-C618-hipotesis, ES-C618-en-teoria, ES-C618-posiblemente, ES-C618-seguramente

\chapter{Supongamos que...}
\label{ch:es-supongamos-que}
\hlchaptermodality{618}

\begin{chapteropening}
\textbf{By the end of this chapter:} \emph{I can imagine a situation and talk about what would probably happen.}

Complete the last lesson of chapter 618: I can imagine a situation and talk about what would probably happen.
\end{chapteropening}

\section[imaginar]{imaginar — to imagine}

\label{lesson:ES-C618-imaginar}



\begin{quote}
Before the new one: say the Spanish for the priority, then the Spanish for the criterion.
\end{quote}

\subsection*{You'll want to know: imaginar}
\textbf{imaginar} — \textquotedblleft{}to imagine\textquotedblright{}.

\begin{quote}
Imagina que ganas la lotería.
\end{quote}

\begin{quote}
\emph{Imagine that you win the lottery.}
\end{quote}

\emph{Me imagino que} — I suppose that.

\begin{cousinweb}
\emph{imaginar} is a learned borrowing from Latin \emph{imaginari} 'to picture to oneself', built on \emph{imago} 'image' (EtymDB).
\end{cousinweb}

\begin{grammarlens}[title={What you've built}]
One more word for telling, explaining and describing at B1.
\end{grammarlens}

\subsection*{Your turn}
\begin{itemize}
\raggedright
  \item \emph{Say it:} \emph{imaginar}
  \item \emph{Say it:} \emph{imaginar}, once more
  \item \emph{Say it:} a sentence of your own with \emph{imaginar}
  \item \emph{Recall:} say the Spanish for the priority, then the Spanish for the criterion, then say \emph{imaginar} again
\end{itemize}

\begin{tcolorbox}[breakable,colback=teal!4,colframe=teal!35!black,title={Before you move on}]
What does imaginar mean? (\textquotedblleft{}To imagine\textquotedblright{}.) Say it once more.
\end{tcolorbox}
\section[la hipótesis]{la hipótesis — the hypothesis}

\label{lesson:ES-C618-hipotesis}



\begin{quote}
Before the new one: say the Spanish for the criterion, then the Spanish for to imagine.
\end{quote}

\subsection*{You'll want to know: la hipótesis}
\textbf{la hipótesis} — \textquotedblleft{}the hypothesis\textquotedblright{}.

\begin{quote}
Es solo una hipótesis.
\end{quote}

\begin{quote}
\emph{It is only a hypothesis.}
\end{quote}

The same form in the plural: \emph{las hipótesis}.

\begin{cousinweb}
\emph{hipótesis} is a learned borrowing from Late Latin \emph{hypothesis}, taken from Greek \emph{hypothesis} 'a base, a supposition', from \emph{hypo} 'under' and \emph{thesis} 'a placing' (EtymDB).
\end{cousinweb}

\begin{grammarlens}[title={What you've built}]
One more word for telling, explaining and describing at B1.
\end{grammarlens}

\subsection*{Your turn}
\begin{itemize}
\raggedright
  \item \emph{Say it:} \emph{la hipótesis}
  \item \emph{Say it:} \emph{la hipótesis}, once more
  \item \emph{Say it:} a sentence of your own with \emph{la hipótesis}
  \item \emph{Recall:} say the Spanish for the criterion, then the Spanish for to imagine, then say \emph{la hipótesis} again
\end{itemize}

\begin{tcolorbox}[breakable,colback=teal!4,colframe=teal!35!black,title={Before you move on}]
What does la hipótesis mean? (\textquotedblleft{}The hypothesis\textquotedblright{}.) Say it once more.
\end{tcolorbox}
\section[en teoría]{en teoría — in theory}

\label{lesson:ES-C618-en-teoria}



\begin{quote}
Before the new one: say the Spanish for to imagine, then the Spanish for the hypothesis.
\end{quote}

\subsection*{You'll want to know: en teoría}
\textbf{en teoría} — \textquotedblleft{}in theory\textquotedblright{}.

\begin{quote}
En teoría, el tren llega a las ocho.
\end{quote}

\begin{quote}
\emph{In theory, the train arrives at eight.}
\end{quote}

It often goes with \emph{en la práctica}.

\begin{cousinweb}
\emph{en teoría} is built on \emph{teoría}, from Late Latin \emph{theoria}, taken from Greek \emph{theōria} 'a looking at, contemplation' (EtymDB).
\end{cousinweb}

\begin{grammarlens}[title={What you've built}]
One more word for telling, explaining and describing at B1.
\end{grammarlens}

\subsection*{Your turn}
\begin{itemize}
\raggedright
  \item \emph{Say it:} \emph{en teoría}
  \item \emph{Say it:} \emph{en teoría}, once more
  \item \emph{Say it:} a sentence of your own with \emph{en teoría}
  \item \emph{Recall:} say the Spanish for to imagine, then the Spanish for the hypothesis, then say \emph{en teoría} again
\end{itemize}

\begin{tcolorbox}[breakable,colback=teal!4,colframe=teal!35!black,title={Before you move on}]
What does en teoría mean? (\textquotedblleft{}In theory\textquotedblright{}.) Say it once more.
\end{tcolorbox}
\section[posiblemente]{posiblemente — possibly}

\label{lesson:ES-C618-posiblemente}



\begin{quote}
Before the new one: say the Spanish for the hypothesis, then the Spanish for in theory.
\end{quote}

\subsection*{You'll want to know: posiblemente}
\textbf{posiblemente} — \textquotedblleft{}possibly\textquotedblright{}.

\begin{quote}
Posiblemente llegue mañana.
\end{quote}

\begin{quote}
\emph{He will possibly arrive tomorrow.}
\end{quote}

Often with the subjunctive.

\begin{cousinweb}
\emph{posiblemente} is built on \emph{posible}, a learned borrowing from Latin \emph{possibilis}, built on \emph{posse} 'to be able', with the ending \emph{-mente} (EtymDB).
\end{cousinweb}

\begin{grammarlens}[title={What you've built}]
One more word for telling, explaining and describing at B1.
\end{grammarlens}

\subsection*{Your turn}
\begin{itemize}
\raggedright
  \item \emph{Say it:} \emph{posiblemente}
  \item \emph{Say it:} \emph{posiblemente}, once more
  \item \emph{Say it:} a sentence of your own with \emph{posiblemente}
  \item \emph{Recall:} say the Spanish for the hypothesis, then the Spanish for in theory, then say \emph{posiblemente} again
\end{itemize}

\begin{tcolorbox}[breakable,colback=teal!4,colframe=teal!35!black,title={Before you move on}]
What does posiblemente mean? (\textquotedblleft{}Possibly\textquotedblright{}.) Say it once more.
\end{tcolorbox}
\section[seguramente]{seguramente — most likely, probably}

\label{lesson:ES-C618-seguramente}



\begin{quote}
Before the new one: say the Spanish for in theory, then the Spanish for possibly.
\end{quote}

\subsection*{You'll want to know: seguramente}
\textbf{seguramente} — \textquotedblleft{}most likely, probably\textquotedblright{}.

\begin{quote}
Seguramente está en casa.
\end{quote}

\begin{quote}
\emph{She is probably at home.}
\end{quote}

It means probably, not surely: \emph{seguramente} leaves room for doubt.

\begin{cousinweb}
\emph{seguramente} is built on \emph{seguro} and \emph{-mente}; \emph{seguro} is inherited from Latin \emph{securus} 'free from care' (EtymDB).
\end{cousinweb}

\begin{grammarlens}[title={What you've built}]
One more word for telling, explaining and describing at B1.
\end{grammarlens}

\subsection*{Your turn}
\begin{itemize}
\raggedright
  \item \emph{Say it:} \emph{seguramente}
  \item \emph{Say it:} \emph{seguramente}, once more
  \item \emph{Say it:} a sentence of your own with \emph{seguramente}
  \item \emph{Recall:} say the Spanish for in theory, then the Spanish for possibly, then say \emph{seguramente} again
\end{itemize}

\begin{tcolorbox}[breakable,colback=teal!4,colframe=teal!35!black,title={Before you move on}]
What does seguramente mean? (\textquotedblleft{}Most likely, probably\textquotedblright{}.) Say it once more.
\end{tcolorbox}
